\section{Results}
\label{sec:results}

\subsection{Clustering}

\subsubsection{Principal Component Analysis}

As specified in Section \ref{sec:data_trbs}, we performed a PCA to reduce the dimensionality of our dataset. 
We set an explained variance threshold target of 80\%, which corresponds to 14 principal components.
This thresholds returns a number of principal components 
Please refer to Table \ref{tab:pca} for further details about the principal components.

\begin{table}[htbp]
  \footnotesize
  \caption{Principal Component details}
  \label{tab:pca}
  \begin{tabular}{rrll}
\toprule
PC & EV* & Most Impacting Features & Description \\
\midrule
\multirow{2}{*}{1} & \multirow{2}{*}{0.240} & Drib. per touch ($\uparrow$), OREB contest\% ($\downarrow$), Pass. rec. ($\uparrow$), & Perimeter \\
& & Elbow touches ($\downarrow$), 2FG contest ($\downarrow$), Screen AST ($\downarrow$), FG3A ($\uparrow$) & ball handler \\
\midrule
\multirow{2}{*}{2} & \multirow{2}{*}{0.144} & Post FG\% ($\uparrow$), Min. ($\uparrow$), Elbow FG\% ($\uparrow$), Touches ($\uparrow$), Pass. rec. ($\uparrow$), & High usage \\
& & Pass. made ($\uparrow$), AST ($\uparrow$), Paint FG\% ($\uparrow$), USG\% ($\uparrow$) & inside player \\
\midrule
\multirow{2}{*}{3} & \multirow{2}{*}{0.088} & Elbow FG\% ($\uparrow$), Paint FG\% ($\uparrow$), Post FG\% ($\uparrow$), Min. ($\uparrow$), Pass. made ($\downarrow$), & Efficient shooter \\
& & Touches ($\downarrow$), Pass. rec. ($\downarrow$), FG3\% ($\uparrow$), Potential AST ($\downarrow$), FG3A ($\uparrow$) & with low poss. \\
\midrule
\multirow{2}{*}{4} & \multirow{2}{*}{0.058} & Post FG\% ($\uparrow$), Elbow FG\% ($\downarrow$), Paint FG\% ($\downarrow$), 2FG contest ($\downarrow$), & Poor defensive \\
& & Deflections ($\downarrow$), Def Boxouts ($\downarrow$), Post PTS ($\uparrow$), Screen AST ($\downarrow$), Min. ($\uparrow$) & post shoot \\ 
\midrule
5 & 0.045 & 2FG contest ($\uparrow$), Post FG\% ($\uparrow$), Def Boxouts ($\uparrow$), Deflections ($\uparrow$) & Inside defense \\ 
\midrule
6 & 0.040 & Post FG\% ($\uparrow$), Min. ($\downarrow$), Elbow FG\% ($\uparrow$), FG3A ($\downarrow$), FG3\% ($\downarrow$), Drive FG\% ($\downarrow$) & Inside shoot \\ 
\midrule
7 & 0.031 & Elbow FG\% ($\uparrow$), Paint FG\% ($\downarrow$), OREB contest\% ($\downarrow$), Drive FG\% ($\downarrow$) & Low hustle shoot \\ 
\midrule
8 & 0.029 & Drive FG\% ($\uparrow$), OREB contest\% ($\downarrow$), Pull-up FG\% ($\uparrow$), FG3A ($\downarrow$), FG2\% ($\uparrow$) & Shot creation \\ 
\midrule
9 & 0.027 & Paint FG\% ($\uparrow$), Elbow FG\% ($\downarrow$), Deflections ($\downarrow$), Catch \& shoot \% ($\uparrow$), FG3\% ($\uparrow$) & Catch \& shoot \\ 
\midrule
10 & 0.024 & OREB contest\% ($\uparrow$), Paint FG\% ($\downarrow$), Drive FG\% ($\uparrow$), Pass. made ($\uparrow$) & Outside shoot \& hustle \\ 
\midrule
11 & 0.023 & Min. ($\uparrow$), USG\% ($\downarrow$), Paint AST ($\uparrow$), FG3A ($\downarrow$), FG2A ($\downarrow$), Deflections ($\uparrow$) & High min., low resp. defense \\ 
\midrule
12 & 0.019 & Pull-up FG\% ($\uparrow$), Deflections ($\uparrow$), Def Rim FG\% ($\uparrow$), FG3A ($\downarrow$), STL ($\uparrow$) & Polyvalent defense \\ 
\midrule
13 & 0.018 & Def Rim FG\% ($\uparrow$), Deflections ($\downarrow$), Post AST ($\downarrow$), Drive FG\% ($\downarrow$), Min. ($\uparrow$) & Rim defense \\ 
\midrule
14 & 0.018 & Def Rim FG\% ($\uparrow$), Pull-up FG\% ($\downarrow$), Drive FG\% ($\uparrow$), Def Boxouts ($\uparrow$) & Inside defense \& shot creation \\ 
\bottomrule
\end{tabular}

  {\raggedright \small $\uparrow$ (respectively $\downarrow$) indicates feature positively (resp. negatively) correlated with the corresponding PC.}
\end{table}


\subsubsection{Clustering algorithms}

The Calinski-Harabasz indices of the compared method are presented in Figure \ref{fig:cl_ch}.
Please note that the number of clusters if for some algorithms (e.g. K-Means, GMM) an input hyperparameter, but for some others (e.g. Partition - HDBSCAN), it is a consequence of other hyperparameters.
\hl{A clear hierarchy emerges between algorithm : K-Means, GMM and Spherical K-Means consistently achieve top performance, while Agglomerative with Complete and Average linkage emerge as worst performers.
Agglomerative with Ward linkage and Partition HDBSCAN fall below the top three performers, the former appearing more constant (and better in most cluster sizes) than the latter.}
From this analysis, we will use K-Means as partition algorithm for the remaining of the paper.

\ob{Choose nb of clusters}

\begin{figure}[h]
  \includegraphics[width=\textwidth]{clustering_ch.png}  
  \caption{Comparison of clustering algorithm.}
  \label{fig:cl_ch}
\end{figure}

Figure \ref{fig:cl_viz} presents a visual inspection of the selected partition. 

\begin{figure}[h]
  \includegraphics[width=\textwidth]{clustering_inspection.png}  
  \caption{Visualization of the selected partition. The left panel represents players in a 2D space determined by the two first principal components and the right panel a histogram of cluster populations.}
  \label{fig:cl_viz}
\end{figure}


\subsection{Roster Constitution}

\ob{To do}